%------------------------------------------------------------------------------%
\documentclass[a4paper,12pt,fleqn]{article}

\usepackage{calculo_varias_variaveis-1}

\ShowAnswers

%------------------------------------------------------------------------------%
\begin{document}

\makeHeader{CFVVI}{Exame Especial -- Turma 4}{28/07/2025}

% 01
%------------------------------------------------------------------------------%
\exercise{25}
Considerando que a equação
\(
  \sin(x+y) + \sin(y+z) + \sin(x+z) = 0
\)
define $z$ como função de $x$ e $y$.
\begin{enumerate}[label=\alph*)]
  \item Calcule as derivadas \(\D{z}{x}\) e \(\D{z}{y}\)
  \item Encontre os valores de \(\D{z}{x}\) e \(\D{z}{y}\) em \((\pi, \pi, \pi)\)
  \item Construa a aproximação linear da função $z(x, y)$ no ponto \((\pi, \pi)\)
\end{enumerate}
\clearpagequestiononly

\begin{answer}
  \setlength{\jot}{6pt}
  \vspace{\baselineskip}
  \textbf{a)\;}
  Considerando $z=z(x, y)$ e derivando por $x$ os dois lados da equação temos
  \begin{align*}
    \sin(x+y) + \sin(y+z(x, y)) + \sin(x+z(x, y))                                     & = 0 \\
    \D{}{x}\left(\sin(x+y) + \sin(y+z(x, y)) + \sin(x+z(x, y))\right)                 & = 0 \\
    \cos(x+y)\D{}{x}(x+y) + \cos(y+z(x, y))\D{}{x}(y+z(x, y)) + \cos(x+z(x, y))\D{}{x}(x+z(x, y)) & = 0 \\
    \cos(x+y) + \cos(y+z(x, y))\D{z}{x} + \cos(x+z(x, y))\left(1+\D{z}{x}\right)      & = 0 \\
    \cos(x+y) + \cos(y+z(x, y))\D{z}{x} + \cos(x+z(x, y)) + \cos(x+z(x, y))\D{z}{x}   & = 0 \\
    \left(\cos(y+z(x, y))+\cos(x+z(x, y))\right)\D{z}{x}  = -\cos(x+y) - \cos(x+z(x, y))    \\
    \D{z}{x}  = -\frac{\cos(x+y)+\cos(x+z(x, y))}{\cos(y+z(x, y))+\cos(x+z(x, y))}
  \end{align*}

  \clearpage
  Derivando por $y$ os dois lados da equação temos
  \begin{align*}
    \sin(x+y) + \sin(y+z(x, y)) + \sin(x+z(x, y))                                     & = 0 \\
    \D{}{y}\left(\sin(x+y) + \sin(y+z(x, y)) + \sin(x+z(x, y))\right)                 & = 0 \\
    \cos(x+y)\D{}{y}(x+y) + \cos(y+z(x, y))\D{}{y}(y+z(x, y)) + \cos(x+z(x, y))\D{}{y}(x+z(x, y)) & = 0 \\
    \cos(x+y) + \cos(y+z(x, y))\left(1+\D{z}{y}\right) + \cos(x+z(x, y))\D{z}{y}      & = 0 \\
    \cos(x+y) + \cos(y+z(x, y)) + \cos(y+z(x, y))\D{z}{y} + \cos(x+z(x, y))\D{z}{y}   & = 0 \\
    \left(\cos(y+z(x, y)) + \cos(x+z(x, y))\right)\D{z}{y}  = -\cos(x+y)-\cos(y+z(x, y))    \\
    \D{z}{y}  = -\frac{\cos(x+y)+\cos(y+z(x, y))}{\cos(y+z(x, y)) + \cos(x+z(x, y))}
  \end{align*}

  \vspace{\baselineskip}
  \textbf{b)\;}
  Avaliando as derivadas no ponto \((\pi, \pi, \pi)\)
  \begin{align*}
    \D{z}{x}(\pi, \pi)
    & = \left(-\frac{\cos(x+y)+\cos(x+z)}{\cos(y+z)+\cos(x+z)}\right)\evalat{(\pi, \pi, \pi)}{} \\
    & = -\frac{\cos(\pi+\pi)+\cos(\pi+\pi)}{\cos(\pi+\pi)+\cos(\pi+\pi)} \\
    & = -\frac{2\cos(2\pi)}{2\cos(2\pi)}  \\
    & = -\frac{2\times 1}{2\times 1}  = -1
    \\
    \D{z}{y}(\pi, \pi)
    & = \left(-\frac{\cos(x+y)+\cos(y+z)}{\cos(y+z)+\cos(x+z)}\right)\evalat{(\pi, \pi, \pi)}{} \\
    & = -\frac{\cos(\pi+\pi)+\cos(\pi+\pi)}{\cos(\pi+\pi)+\cos(\pi+\pi)} \\
    & = -\frac{2\cos(2\pi)}{2\cos(2\pi)}  \\
    & = -\frac{2\times 1}{2\times 1}  = -1
  \end{align*}

  \vspace{\baselineskip}
  \textbf{c)\;}
  A aproximação linear da função $z(x, y)$ no ponto \((x_0, y_0) = (\pi, \pi)\) é
  \begin{align*}
    L(x, y)
    & = z(x_0, y_0) + (x-x_0)\D{z}{x}(x_0, y_0) + (y-y_0)\D{z}{y}(x_0, y_0) \\
    & = z(\pi, \pi) + (x-\pi)\D{z}{x}(\pi, \pi) + (y-\pi)\D{z}{y}(\pi, \pi) \\
    & = \pi -(x-\pi) - (y-\pi) \\
    & = \pi -x +\pi -y +\pi \\
    & = 3\pi -x -y
  \end{align*}
\end{answer}

% 02
%------------------------------------------------------------------------------%
\exercise{25}
Encontre os valores máximo e mínimo da função
\(
  f(x, y) = x^2 + 4y^2
\)
restrita a região fechada limitada
\(
  x^2 + y^2 \leq 9
\)
\clearpagequestiononly

\begin{answer}
  \setlength\columnsep{3em}
  \setlength{\jot}{5pt}
  \begin{multicols}{2}
    \raggedcolumns
    Como a função é contínua e a região é fechada e limitada sabemos que
    $f$ assume um valor máximo e um valor mínimo na região.

    \vspace{\baselineskip}
    Temos que buscar os pontos críticos no interior e aplicar multiplicadores
    de Lagrange na fronteira.

    \vspace{\baselineskip}
    \textbf{Gradiente} de $f$
    \[
      \D{f}{x}
      = \D{}{x}\left(x^2 + 4y^2\right)
      = 2x
    \]
    \[
      \D{f}{y}
      = \D{}{y}\left(x^2 + 4y^2\right)
      = 8y
    \]

    \vspace{\baselineskip}
    \textbf{Pontos críticos}:
    Impondo \(\nabla f = 0\) temos o ponto interior
    \[
      P_1 = (0, 0)
    \]

    \vspace{\baselineskip}
    \textbf{Multiplicadores de Lagrange}:
    Temos a restrição \(g(x, y) = x^2 + y^2 \leq 9\) e seu gradiente é
    \[
      \D{g}{x}
      = \D{}{x}\left(x^2 + y^2\right)
      = 2x
    \]
    \[
      \D{g}{y}
      = \D{}{y}\left(x^2 + y^2\right)
      = 2y
    \]

    Precisamos resolver o sistema
    \begin{gather*}
      2x = 2\lambda x \\
      8y = 2\lambda y \\
      x^2 + y^2 = 9
    \end{gather*}
    Que pode ser simplificado
    \begin{gather*}
      x = \lambda x \\
      4y = \lambda y \\
      x^2 + y^2 = 9
    \end{gather*}
    Da primeira equação temos que $x=0$ ou $\lambda=1$.

    Se $x=0$ a terceira equação se reduz a $y^2=9$, cujas soluções são
    $y=3$ ou $y=-3$. Substituindo qualquer uma delas na segunda equação
    temos $\lambda = 0$. Encontramos o segundo e terceiro pontos
    \[
      P_2 = (0, -3)
      \qquad
      P_3 = (0, 3)
    \]

    Se $\lambda=1$ a segunda equação se torna $4y=y$ e, portanto, $y=0$.
    Substituindo na terceira equação temos $x^2=9$, cujas soluções são
    $x=-3$ ou $x=3$.  Encontramos o quarto e quinto pontos
    \[
      P_4 = (-3, 0)
      \qquad
      P_5 = (3, 0)
    \]

    \textbf{Avaliando} a função nos pontos encontrados
    \begin{gather*}
      f(0, 0)  =    0^2 + 4\times    0^2 = 0 \\
      f(0, -3) =    0^2 + 4\times (-3)^2 = 36 \\
      f(0, 3)  =    0^2 + 4\times    3^2 = 36 \\
      f(-3, 0) = (-3)^2 + 4\times    0^2 = 9 \\
      f(3, 0)  =    3^2 + 4\times    0^2 = 9
    \end{gather*}

    O valor mínimo de $f$ é 0 e o valor máximo é 36
  \end{multicols}
\end{answer}

% 3
%------------------------------------------------------------------------------%
\exercise{25}
Considerando a função
\(
  f(x, y) = 4 - \sqrt{16-x^2-y^2}
\)
\begin{enumerate}[label=\alph*)]
  \item Determine e esboce o domínio de $f$
  \item Caracterize todas as curvas de nível de $f$ e esboce três delas
  \item Calcule o limite
  \(
    \lim_{(x, y)\to(0, 0)} \frac{x^2+y^2}{f(x, y)}
  \)
  ou mostre que ele não existe
\end{enumerate}
\begin{questiononly}
  \begin{multicols}{2}
    Domínio
    \vspace{-0.8\baselineskip}
    \begin{center}
      \includegraphics{axis_xy}
    \end{center}
    Curvas de nível
    \vspace{-0.8\baselineskip}
    \begin{center}
      \includegraphics{axis_xy}
    \end{center}
  \end{multicols}
\end{questiononly}
\clearpagequestiononly

\begin{answer}
  \begin{multicols}{2}
    Domínio
    \vspace{-0.8\baselineskip}
    \begin{center}
      \includegraphics{T2_dominio}
    \end{center}
    Curvas de nível
    \vspace{-0.8\baselineskip}
    \begin{center}
      \includegraphics{T2_curvas_nivel}
    \end{center}
  \end{multicols}

  \begin{multicols}{2}
    \raggedcolumns
    \setlength{\jot}{5pt}

    \vspace{\baselineskip}
    \textbf{a)\;}
    O domínio de $f$ consiste dos pontos onde é possível avaliar
    a raiz quadrada, isto é, os pontos que satisfazem $16-x^2-y^2\geq 0$,
    ou seja,
    \begin{align*}
      16 -x^2 -y^2 & \geq 0   \\
      -x^2 -y^2    & \geq -16 \\
      x^2 +y^2     & \leq 16 = 4^2
    \end{align*}
    Portanto,
    \[
      D_f = \left\{(x, y)\in\R^2 \st x^2 + y^2 \leq 16 \right\},
    \]
    que corresponde a um disco de raio 4, centrado na origem.

    \vspace{\baselineskip}
    \textbf{b)\;}
    As curvas de nível de $f$ são compostas pelos pontos onde $f(x, y) = c$
    para alguma constante $c$ na imagem de $f$.

    Para determinar a imagem de $f$ observamos que $x^2+y^2$ só pode assumir
    valores no intervalo $[0, 16]$. Portanto, $16 - x^2 - y^2$ está limitado
    ao mesmo intervalo. Consequentemente $\sqrt{16 - x^2 - y^2}$ está em $[0, 4]$.
    Concluímos que os valores de $f$ estão em $[0, 4]$, isto é,
    \(\Im(f) = [0, 4]\)

    Para qualquer $c\in\Im(f)=[0, 4]$, temos a curva de nível
    \begin{align*}
      f(x, y)                 & = c              \\
      4 - \sqrt{16 -x^2 -y^2} & = c              \\
      \sqrt{16 -x^2 -y^2}     & = 4 - c          \\
      16 -x^2 -y^2            & = (4 - c)^2      \\
      -x^2 -y^2               & = (4 - c)^2 -16  \\
      x^2 +y^2                & = 16 - (4 - c)^2
    \end{align*}
    Que corresponde a uma circunferência centrada na origem de raio
    \[
      r = \sqrt{16 - (4 - c)^2}
    \]

    Escolhendo os valores $c=1$, $c=2$ e $c=4$, temos as curvas de nível
    \begin{align*}
      \gamma_1\colon & x^2 + y^2 = 16 - (4 - 1)^2 = 16 - 9 = 7  \\
      \gamma_2\colon & x^2 + y^2 = 16 - (4 - 2)^2 = 16 - 4 = 12 \\
      \gamma_3\colon & x^2 + y^2 = 16 - (4 - 4)^2 = 16 - 0 = 16
    \end{align*}

    \vspace{\baselineskip}
    \textbf{c)\;}
    Calculando o limite
    \setlength{\jot}{12pt}
    \setlength{\mathindent}{0cm}
    \begin{align*}
      L
      & = \lim_{(x, y)\to(0, 0)} \frac{x^2+y^2}{f(x, y)} \\
      & = \lim_{(x, y)\to(0, 0)} \frac{x^2+y^2}{4 - \sqrt{16-x^2-y^2}} \\
      & = \lim_{(x, y)\to(0, 0)} \frac{x^2+y^2}{4 - \sqrt{16-x^2-y^2}} \times \\
      & \hspace{32mm}            \frac{4 + \sqrt{16-x^2-y^2}}{4 + \sqrt{16-x^2-y^2}} \\
      & = \lim_{(x, y)\to(0, 0)} \frac{\left(x^2+y^2\right)\left(4 + \sqrt{16-x^2-y^2}\right)}
                                      {4^2 - \left(\sqrt{16-x^2-y^2}\right)^2} \\
      & = \lim_{(x, y)\to(0, 0)} \frac{\left(x^2+y^2\right)\left(4 + \sqrt{16-x^2-y^2}\right)}
                                      {16 - (16-x^2-y^2)} \\
      & = \lim_{(x, y)\to(0, 0)} \frac{x^2+y^2}{x^2+y^2}\left(4 + \sqrt{16-x^2-y^2}\right) \\
      & = \lim_{(x, y)\to(0, 0)} \left(4 + \sqrt{16-x^2-y^2}\right) \\
      & = 4 + \sqrt{16-0^2-0^2} \\
      & = 4 + 4 = 8
    \end{align*}
  \end{multicols}
\end{answer}

% 04
%------------------------------------------------------------------------------%
\exercise{25}
Dado \(z=\frac{ \left(1-\sqrt{3}\right) + \left(1+\sqrt{3}\right)i}{2+2i} \), calcule
\begin{enumerate}[label=\alph*),itemsep=0pt]
  \item a parte real de $z$,
  \item a parte imaginária de $z$,
  \item o módulo de $z$,
  \item o argumento de $z$
\end{enumerate}
\clearpagequestiononly

\begin{answer}
  Avaliando $z$
  \begin{align*}
    z
    & = \frac{ \left(1-\sqrt{3}\right) + \left(1+\sqrt{3}\right)i}{2+2i} \\[2mm]
    & = \frac{ \left[\left(1-\sqrt{3}\right) + \left(1+\sqrt{3}\right)i\right](2-2i)}{(2+2i)(2-2i)} \\[2mm]
    & = \frac{ \left(1-\sqrt{3}\right)(2-2i) + \left(1+\sqrt{3}\right)(2i+2)}{2^2 + 2^2} \\[2mm]
    & = \frac{ 2-2i-2\sqrt{3}+2\sqrt{3}i + 2i+2+2\sqrt{3}i +2\sqrt{3}}{8} \\[2mm]
    & = \frac{ 4 + 4\sqrt{3}i}{8} \\[2mm]
    & = \frac{1}{2} + \frac{\sqrt{3}}{2}i
  \end{align*}
  Parte real de $z$
  \[
    \Re(z) = \frac{1}{2}
  \]
  Parte imaginária de $z$
  \[
    \Im(z) = \frac{\sqrt{3}}{2}
  \]
  Módulo de $z$
  \[
    \abs{z}
    = \sqrt{\left(\frac{1}{2}\right)^2 + \left(\frac{\sqrt{3}}{2}\right)^2}
    = \sqrt{\frac{1}{4} + \frac{3}{4}}
    = \sqrt{\frac{1+3}{4}}
    = \sqrt{1}
    = 1
  \]
  Argumento de $z$, \(\varphi=\arg(z)\)
  \[
    \sin(\varphi)
    = \frac{\Im(z)}{\abs{z}}
    = \frac{\sqrt{3}}{2}
    \qquad\qquad
    \cos(\varphi)
    = \frac{\Re(z)}{\abs{z}}
    = \frac{1}{2}
  \]
  \[
    \arg(z)
    = \ang{60}
    = \frac{\pi}{3}
  \]
\end{answer}

%------------------------------------------------------------------------------%
\end{document}
%------------------------------------------------------------------------------%
